problem:  
Let \( (M^{2n}, J, g_0) \) be a **compact complex manifold** admitting a **strong Kähler with torsion (SKT)** metric \( g_0 \), with fundamental form \( \omega_0 \) satisfying \( \partial\bar{\partial}\omega_0=0 \).  
Assume that \( M \) has **trivial canonical bundle**, i.e., there exists a nowhere‑vanishing holomorphic volume form \( \Omega \) with \( \bar{\partial}\Omega=0 \).

Consider the **normalized pluriclosed flow**  
\[
\frac{\partial \omega_t}{\partial t} = -(\rho^B)^{1,1} + \lambda_t \omega_t,
\]
where \( \rho^B \) denotes the Bismut–Ricci form and the normalization coefficient  
\[
\lambda_t = \frac{1}{2n\,\mathrm{Vol}(M,\omega_t)}\int_M \operatorname{tr}_{\omega_t}\big((\rho^B)^{1,1}\big)\,dV_{\omega_t}
\]
ensures that the total volume of \( (M, g_t) \) remains constant in time.  

Suppose further that \( M \) belongs to a structurally well‑controlled class, such as a **complex nilmanifold** or a **complex surface**, ensuring long‑time existence of the pluriclosed flow (e.g., as in known results by Streets–Tian and Ustinovskiy).  

**Question:**  
In these settings, does the normalized pluriclosed flow starting from \( \omega_0 \) converge smoothly (up to diffeomorphism) to a limiting SKT metric \( \omega_\infty \) satisfying  
\[
(\rho^B)^{1,1}=0, \quad \text{and}\quad \nabla^B\Omega=0?
\]
If not always, can one identify analytic or cohomological criteria intrinsic to \( (M,J,[\omega_0]) \) that determine whether such convergence to a **Bismut–Ricci‑flat structure with holonomy \(\subseteq SU(n)\)** can occur?

Why is it a "good" problem:  
This refined problem focuses the challenge on classes of SKT manifolds where pluriclosed flow is known to exist and evolve smoothly—thereby making the convergence conjecture both well‑posed and approachable with current tools. It asks for a **characterization of geometric and analytic conditions** under which pluriclosed flow produces a Bismut–Ricci‑flat, holonomy‑reduced structure.  

The question aims squarely at a non‑Kähler analogue of the **Calabi–Yau theorem** within the SKT framework and at understanding the possible dynamical mechanisms enforcing holonomy reduction under geometric evolution. It links together three major threads:
- nonlinear PDE methods (pluriclosed flow and a priori estimates),
- Hermitian holonomy theory (parallel \((n,0)\) forms and \(SU(n)\)-structures),
- and the cohomological landscape (Bott–Chern and Aeppli classes).

It is **simple to state**, depending only on standard geometric flows and holonomy, but it requires deep insight across complex geometry and geometric analysis to progress. Resolving it would clarify whether pluriclosed flow has the canonical–metric–producing power in the SKT world that Ricci flow enjoys in Kähler geometry, making it an archetypal “good” research problem: conceptually clear, cross‑disciplinary, and foundational for understanding the structure of non‑Kähler geometry.